%================================================================
\chapter{Entornos virtuales de Python y automatización del flujo de trabajo}
\label{ch:python-environments}
%================================================================

Antes de escribir cualquier script, necesitamos comprender \emph{qué} van a
gestionar esos scripts: los entornos virtuales de Python.

\textbf{El problema: conflictos entre paquetes}.
Imagine que tiene dos proyectos en su computadora. El proyecto~A usa la
versión~1.5 de una biblioteca llamada \texttt{numpy}; el proyecto~B requiere
la versión~2.0. Si instala ambas versiones de forma global, entrarán en
conflicto. Gana la última que se haya instalado, y uno de sus proyectos
podría dejar de funcionar.

\textbf{La solución: entornos virtuales}.
Un \textbf{entorno virtual} (o \emph{venv}) es una copia aislada de Python
con su propio conjunto de paquetes instalados, completamente separada de
cualquier otro entorno en su máquina. Piénselo como una burbuja autocontenida
alrededor de su proyecto.

\begin{itemize}
  \item Cada proyecto obtiene su propio entorno.
  \item Los paquetes instalados en un entorno no pueden interferir con otro.
  \item Puede recrear un entorno de manera exacta a partir de una lista de paquetes.
  \item Eliminar un entorno es seguro y limpio.
\end{itemize}

Cubriremos dos estrategias para gestionar entornos: el enfoque tradicional
mediante \texttt{pip} y una tendencia emergente llamada \texttt{uv}.

%---------------------------------------------------------
\section{Gestión tradicional de paquetes: \texttt{pip}}
\label{sec:venvs}
%---------------------------------------------------------
\subsubsection{Windows (CMD o PowerShell)}

\begin{lstlisting}[style=windowsStyle, caption={Creación de un venv en Windows},escapeinside={(*@}{@*)}]
:: Navegue a la carpeta de su proyecto
cd C:\Projects\(*@\ProjectName@*)

:: Cree el entorno (Python debe estar instalado globalmente)
python -m venv C:\penv\(*@\ProjectName@*)
\end{lstlisting}

\subsubsection{Linux / WSL}

\begin{lstlisting}[style=linuxStyle, caption={Creación de un venv en Linux/WSL},escapeinside={(*@}{@*)}]
# Navegue a la carpeta de su proyecto
cd ~/(*@\ProjectName@*)

# Cree el entorno
python3 -m venv ~/penv/(*@\ProjectName@*)
\end{lstlisting}

En ambos casos, \texttt{C:\textbackslash penv\textbackslash \ProjectName{}} o
\texttt{\textasciitilde/penv/\ProjectName{}} es donde se almacenará el entorno.
Mantener todos los entornos en una única carpeta (\texttt{penv} significa
«Python environments», es decir, «entornos de Python») facilita encontrarlos
y gestionarlos.

\subsection{Activación de un entorno virtual}

Activar un entorno modifica la sesión actual de su intérprete de comandos
de modo que \texttt{python} y \texttt{pip} hagan referencia a las copias
aisladas dentro del venv, y no a las globales.

\subsubsection{Windows}

\begin{lstlisting}[style=windowsStyle, caption={Activación de un venv en Windows},escapeinside={(*@}{@*)}]
C:\penv\(*@\ProjectName@*)\Scripts\activate.bat
\end{lstlisting}

Tras la activación, su indicación cambiará para mostrar el nombre del entorno:
\begin{lstlisting}[style=windowsStyle,escapeinside={(*@}{@*)}]
((*@\ProjectName@*)) C:\Projects\(*@\ProjectName@*)>
\end{lstlisting}

\subsubsection{Linux / WSL}

\begin{lstlisting}[style=linuxStyle, caption={Activación de un venv en Linux/WSL},escapeinside={(*@}{@*)}]
source ~/penv/(*@\ProjectName@*)/bin/activate
\end{lstlisting}

Tras la activación:
\begin{lstlisting}[style=linuxStyle,escapeinside={(*@}{@*)}]
((*@\ProjectName@*)) drg@machine:~/(*@\ProjectName@*)$
\end{lstlisting}

\begin{warningbox}
Observe la palabra clave \texttt{source} en el comando de Linux. Esto es
crítico y se explica en detalle en la Sección~\ref{sec:source-vs-execute}.
\end{warningbox}

\subsection{Instalación de paquetes}

Una vez que el entorno está activo, use \texttt{pip} para instalar paquetes.
La sintaxis es idéntica en ambas plataformas:

\begin{lstlisting}[style=inlineStyle, caption={Instalación de paquetes con pip},escapeinside={(*@}{@*)}]
pip install numpy pandas matplotlib openai langchain
\end{lstlisting}

Para guardar una lista de todos los paquetes instalados (de modo que usted o
un colega puedan recrear el entorno más adelante):

\begin{lstlisting}[style=inlineStyle,escapeinside={(*@}{@*)}]
pip freeze > requirements.txt
\end{lstlisting}

Para recrear un entorno a partir de dicha lista:

\begin{lstlisting}[style=inlineStyle,escapeinside={(*@}{@*)}]
pip install -r requirements.txt
\end{lstlisting}

\subsection{Desactivación}

Para salir del entorno y regresar al intérprete de comandos global (idéntico
en ambas plataformas):

\begin{lstlisting}[style=inlineStyle,escapeinside={(*@}{@*)}]
deactivate
\end{lstlisting}

\subsection{Instalación de las bibliotecas requeridas por \ProjectName{}}
\label{sec:project-libraries}

Con el entorno de \ProjectName{} activo, instale las bibliotecas necesarias
para el reto de datos:

\begin{verbatim}
pip install numpy pandas matplotlib scipy jupyter
\end{verbatim}

Estas cinco bibliotecas son suficientes para todo el análisis del reto de
datos. Las bibliotecas adicionales requeridas para el capítulo sobre modelos
grandes de lenguaje (incluyendo \texttt{openai}, \texttt{anthropic} y
\texttt{langchain}) se instalan por separado en el
Capítulo~\ref{ch:ollama-infrastructure}.

\vs
\noindent
Para verificar la instalación:

\begin{verbatim}
python -c "import numpy, pandas, matplotlib, scipy, jupyter; print('OK')"
\end{verbatim}

\vs
\noindent
\textbf{Sabrá que su entorno principal está completamente funcional cuando:}
\begin{itemize}
    \item \texttt{python --version} devuelve un número de versión.
    \item \texttt{code .} abre Visual Studio Code desde una terminal.
    \item \texttt{git --version} devuelve un número de versión.
    \item Las cinco bibliotecas requeridas se importan sin errores (véase arriba).
\end{itemize}
Los programas agente de modelos grandes de lenguaje (\texttt{agentGroq.py},
\texttt{agentGPT.py}, \texttt{agentClaude.py}) se encuentran en
\texttt{resources/unit\_7/FOO/} dentro del repositorio de \ProjectName{}.
Ejecutarlos requiere claves de API configuradas en el
Capítulo~\ref{ch:ollama-infrastructure}.
\vs

% ============================================================
\section{Gestión moderna de paquetes: \texttt{uv}}
\label{sec:uv}
% ============================================================

\subsection{Por qué \texttt{uv} es relevante hoy}

Las herramientas presentadas en la sección anterior (\texttt{python -m venv}
para la creación de entornos y \texttt{pip} para la instalación de paquetes)
son la base establecida y estable de la gestión de paquetes de Python. Son
mantenidas por la Python Packaging Authority (PyPA) bajo el amparo de la
Python Software Foundation (PSF), una organización sin ánimo de lucro con
una estructura de gobernanza de largo plazo y sin dependencia comercial.
Cada instalación de Python incluye \texttt{pip} por defecto, y saber cómo
usarlo sigue siendo una habilidad esencial independientemente de qué otras
herramientas adopte.

Ha surgido una herramienta más reciente llamada \texttt{uv} que está ganando
adopción rápidamente en entornos profesionales. Comprender ambas
herramientas es importante precisamente porque el panorama está en
transición.

\begin{infobox}
\textbf{Una nota sobre la gobernanza.} \texttt{pip} es mantenido por la
Python Packaging Authority bajo la Python Software Foundation, una
organización sin ánimo de lucro con un modelo de gobernanza comunitaria y
sin dependencia comercial. \texttt{uv} es una herramienta de código abierto
desarrollada y mantenida por Astral, una empresa comercial financiada con
capital de riesgo. Astral ha sido un excelente administrador de la
herramienta y ha contribuido de manera significativa al ecosistema más
amplio de empaquetado de Python. No obstante, su trayectoria a largo plazo
depende de los resultados del negocio de una manera en que los proyectos
respaldados por fundaciones no dependen. Hasta que el ecosistema se
estabilice en torno a un estándar comunitario claro, conocer tanto
\texttt{pip} como \texttt{uv} es la postura profesional prudente.
\end{infobox}

\subsection{Qué ofrece \texttt{uv}}

\texttt{uv} es un reemplazo directo para \texttt{pip} y
\texttt{python -m venv}, reescrito en Rust por motivos de rendimiento. Sus
principales ventajas sobre la cadena de herramientas tradicional son:

\begin{itemize}
  \item \textbf{Velocidad.} La resolución e instalación de paquetes son
        de 10 a 100 veces más rápidas que \texttt{pip} debido a una
        implementación en Rust y a una descarga paralela agresiva.
  \item \textbf{Una caché central.} Los paquetes descargados se almacenan
        una sola vez en una caché global y se reutilizan entre entornos.
        En el mismo sistema de archivos, \texttt{uv} usa enlaces físicos
        (hardlinks) de modo que múltiples entornos que comparten un
        paquete no consumen espacio de disco adicional más allá de la
        primera copia.
  \item \textbf{Archivos de bloqueo.} \texttt{uv} produce un archivo
        \texttt{uv.lock} que registra el grafo completo y resuelto de
        dependencias con hashes criptográficos. Esta es una mejora
        significativa respecto a \texttt{pip freeze > requirements.txt},
        que solo registra nombres y versiones de los paquetes instalados
        sin distinguir dependencias directas de las transitivas, y sin
        verificación por hash.
  \item \textbf{Integración con \texttt{pyproject.toml}.} \texttt{uv}
        trabaja de forma nativa con \texttt{pyproject.toml}, el estándar
        moderno para declarar dependencias de un proyecto, reemplazando
        el flujo de trabajo más antiguo basado en \texttt{requirements.txt}.
  \item \textbf{Gestión de versiones de Python.} \texttt{uv} puede
        descargar y gestionar múltiples versiones de Python de forma
        independiente, eliminando la necesidad de herramientas separadas
        como \texttt{pyenv}.
\end{itemize}

\subsection{Instalación de \texttt{uv}}

\subsubsection{Windows (PowerShell)}

\begin{lstlisting}[style=windowsStyle, caption={Instalación de \texttt{uv} en Windows},escapeinside={(*@}{@*)}]
powershell -ExecutionPolicy ByPass -c "irm https://astral.sh/uv/install.ps1 | iex"
\end{lstlisting}

\texttt{uv} se instala en
\texttt{\%USERPROFILE\%\textbackslash.local\textbackslash bin} y se añade
automáticamente al \texttt{PATH} de su usuario. Cierre y reabra su terminal,
y luego verifique:

\begin{lstlisting}[style=windowsStyle,escapeinside={(*@}{@*)}]
uv --version
\end{lstlisting}

\subsubsection{Linux / WSL}

\begin{lstlisting}[style=linuxStyle, caption={Instalación de \texttt{uv} en Linux/WSL},escapeinside={(*@}{@*)}]
curl -LsSf https://astral.sh/uv/install.sh | sh
\end{lstlisting}

\texttt{uv} se instala en \texttt{\textasciitilde/.local/bin}. Recargue la
configuración de su intérprete de comandos y verifique:

\begin{lstlisting}[style=linuxStyle,escapeinside={(*@}{@*)}]
source ~/.bashrc
uv --version
\end{lstlisting}

\subsection{Creación y gestión de entornos con \texttt{uv}}

La estructura de entorno producida por \texttt{uv venv} es idéntica a la de
un entorno estándar creado con \texttt{python -m venv}. Se encuentran
presentes los mismos scripts \texttt{Scripts\textbackslash activate.bat}
(Windows) y \texttt{bin/activate} (Linux), y la activación no cambia.

\subsubsection{Creación de un entorno}

\begin{lstlisting}[style=windowsStyle, caption={Creación de un entorno gestionado por uv (Windows)},escapeinside={(*@}{@*)}]
uv venv C:\uvenv\(*@\ProjectName@*) --python 3.12
\end{lstlisting}

\begin{lstlisting}[style=linuxStyle, caption={Creación de un entorno gestionado por uv (Linux/WSL)},escapeinside={(*@}{@*)}]
uv venv ~/uvenv/(*@\ProjectName@*) --python 3.12
\end{lstlisting}

La bandera \texttt{--python} especifica la versión de Python. Si esa versión
no está ya instalada, \texttt{uv} la descargará e instalará automáticamente.

\subsubsection{Activación de un entorno}

La activación es idéntica a la de un venv estándar:

\begin{lstlisting}[style=windowsStyle,escapeinside={(*@}{@*)}]
C:\uvenv\(*@\ProjectName@*)\Scripts\activate.bat
\end{lstlisting}

\begin{lstlisting}[style=linuxStyle,escapeinside={(*@}{@*)}]
source ~/uvenv/(*@\ProjectName@*)/bin/activate
\end{lstlisting}

\subsubsection{Instalación de paquetes}

Dentro de un entorno activado, \texttt{uv pip install} reemplaza a
\texttt{pip install}:

\begin{lstlisting}[style=inlineStyle, caption={Instalación de paquetes con uv},escapeinside={(*@}{@*)}]
uv pip install numpy pandas matplotlib
\end{lstlisting}

\subsubsection{Declaración de dependencias con \texttt{pyproject.toml}}

Para proyectos que se compartirán o reproducirán, \texttt{uv} funciona mejor
con un archivo \texttt{pyproject.toml} en la raíz del proyecto que declare
las dependencias de forma explícita:

\begin{lstlisting}[style=inlineStyle, caption={\texttt{pyproject.toml} mínimo},escapeinside={(*@}{@*)}]
[project]
name = "my-project"
version = "0.1.0"
requires-python = ">=3.12"
dependencies = [
    "numpy>=1.26",
    "pandas>=2.2",
    "matplotlib>=3.9",
]
\end{lstlisting}

Con este archivo presente, un único comando instala todas las dependencias
y genera el archivo de bloqueo:

\begin{lstlisting}[style=inlineStyle,escapeinside={(*@}{@*)}]
uv sync
\end{lstlisting}

\texttt{uv sync} lee \texttt{pyproject.toml}, resuelve el grafo completo de
dependencias, instala todos los paquetes y escribe \texttt{uv.lock}.
Ejecutar \texttt{uv sync} nuevamente en otra máquina, o después de un
\texttt{git pull} que haya añadido nuevas dependencias, reproduce el
entorno de forma exacta.

\begin{infobox}
Haga commit de \texttt{pyproject.toml} y \texttt{uv.lock} en su repositorio.
No haga commit de la carpeta del entorno en sí. Añádala a
\texttt{.gitignore}.
\end{infobox}

\newpage
\subsection{\texttt{requirements.txt} frente a \texttt{uv.lock}}

\begin{table}[h]
\centering
\caption{Comparación entre \texttt{requirements.txt} y \texttt{uv.lock}}
\label{tab:req-vs-lock}
\begin{tabular}{@{}p{0.28\textwidth}p{0.30\textwidth}p{0.30\textwidth}@{}}
\toprule
\textbf{Propiedad} & \textbf{\texttt{requirements.txt}} & \textbf{\texttt{uv.lock}} \\
\midrule
Generado por
  & \texttt{pip freeze}
  & \texttt{uv sync} / \texttt{uv lock} \\[4pt]
Registra
  & Solo paquetes instalados
  & Grafo completo de dependencias \\[4pt]
Distingue dependencias directas y transitivas
  & No
  & Sí \\[4pt]
Verificación por hash
  & No
  & Sí \\[4pt]
Reproducibilidad entre plataformas
  & Parcial
  & Sí \\[4pt]
Debe incluirse en git
  & Sí
  & Sí \\
\bottomrule
\end{tabular}
\end{table}

% ============================================================
\section{Nivel 1: archivos de script}
\label{sec:level1}
% ============================================================

Ahora que comprendemos los entornos virtuales, podemos escribir scripts para
activarlos automáticamente. El objetivo es crear dos scripts:

\begin{enumerate}
  \item \textbf{\texttt{ActivateEnv}}: un script reutilizable que acepta
        dos parámetros (nombre del entorno y directorio del proyecto),
        activa el entorno y abre VS Code.
  \item \textbf{\texttt{\ProjectName}}: un script específico del proyecto,
        de una sola línea, que invoca a \texttt{ActivateEnv} con los
        parámetros correctos para el proyecto \ProjectName{}.
\end{enumerate}

\subsection{La solución de Windows: archivos por lotes}

\subsubsection{ActivateEnv.bat}

La versión de Windows, \texttt{ActivateEnv.bat}, se almacena en
\texttt{C:\textbackslash penv} y se coloca en el \texttt{PATH} del sistema.
Acepta dos parámetros: \texttt{\%1} (nombre del entorno) y \texttt{\%2}
(directorio del proyecto).

\begin{lstlisting}[style=windowsStyle, basicstyle=\footnotesize\ttfamily, caption={\texttt{ActivateEnv.bat}: lanzador de entorno reutilizable (Windows)},escapeinside={(*@}{@*)}]
@echo off
setlocal

:: Lea los parámetros
set "ENV_NAME=%~1"
set "BASE_DIR=%~2"
echo Iniciando la activación del entorno '\%ENV_NAME%'...

:: Valide la entrada
if "%ENV_NAME%"=="" (
    echo [ERROR] No se especificó el nombre del entorno.
    goto :eof
)
if "%BASE_DIR%"=="" (
    echo [ERROR] No se especificó el directorio base.
    goto :eof
)

:: Verifique si existe el directorio base
if not exist "%BASE_DIR%" (
    echo [ERROR] Directorio no encontrado: %BASE_DIR%
    goto :eof
)

:: Cambie al directorio de trabajo
cd /d "%BASE_DIR%"

:: Defina la ruta del script de activación
set "VENV_ROOT=C:\penv"
set "VENV_ACTIVATE=%VENV_ROOT%\%ENV_NAME%\Scripts\activate.bat"

:: Active el entorno
if exist "%VENV_ACTIVATE%" (
    echo Activando...
    call "%VENV_ACTIVATE%"
    echo Activación completa.
) else (
    echo [ERROR] No encontrado: %VENV_ACTIVATE%
    goto :eof
)

:: Abra VS Code en el directorio del proyecto
echo Iniciando VS Code...
code .

endlocal
\end{lstlisting}


Elementos clave específicos de Windows:
\begin{itemize}
  \item \texttt{@echo off} suprime el eco de cada comando en la pantalla.
  \item \texttt{setlocal} asegura que los cambios de variables sean locales
        a este script.
  \item \texttt{\%\textasciitilde1} elimina las comillas que rodean al
        primer argumento.
  \item \texttt{call} se usa para invocar otro archivo por lotes y devolver
        el control después.
  \item \texttt{goto :eof} salta al final del archivo (una salida limpia).
\end{itemize}

\subsubsection{\ProjectName{}.bat}

Este archivo reside en el directorio del proyecto \ProjectName{}. Simplemente
invoca a \texttt{ActivateEnv.bat} con los argumentos correctos:

\begin{lstlisting}[style=windowsStyle, caption={\texttt{\ProjectName{}.bat}: lanzador del proyecto (Windows)},escapeinside={(*@}{@*)}]
ActivateEnv (*@\ProjectName@*) "C:\Projects\(*@\ProjectName@*)"
\end{lstlisting}

Debido a que \texttt{C:\textbackslash penv} está en el \texttt{PATH} del
sistema, Windows encuentra \texttt{ActivateEnv.bat} automáticamente. Puede
ejecutar \texttt{\ProjectName{}.bat} haciendo doble clic en él o escribiendo
\texttt{\ProjectName{}} en cualquier ventana de CMD.

\subsection{La solución de Linux: scripts de intérprete de comandos}
\label{subsec:linux-scripts}

\subsubsection{Comprensión de la sintaxis de los scripts de intérprete de comandos}

Un script de intérprete de comandos de Linux es un archivo de texto plano que
contiene comandos de Bash. Cada script debe comenzar con una \textbf{línea
shebang}, una primera línea especial que indica al sistema qué intérprete
usar:

\begin{lstlisting}[style=linuxStyle,escapeinside={(*@}{@*)}]
#!/bin/bash
\end{lstlisting}

Los parámetros pasados a un script se acceden como \texttt{\$1}, \texttt{\$2},
\texttt{\$3}, etc. El nombre propio del script es \texttt{\$0}.

\subsubsection{ActivateEnv.sh}

\begin{lstlisting}[style=linuxStyle, basicstyle=\footnotesize\ttfamily, caption={\texttt{ActivateEnv.sh}: lanzador de entorno reutilizable (Linux/WSL)},escapeinside={(*@}{@*)}]
#!/bin/bash
# ActivateEnv.sh
# Uso: source ~/penv/ActivateEnv.sh <ENV_NAME> <PROJECT_DIR>
# Nota: Debe ser *cargado con source*, no ejecutado. Véase la explicación abajo.

ENV_NAME="$1"
BASE_DIR="$2"

echo "Iniciando la activación del entorno '${ENV_NAME}'..."

# Valide la entrada
if [ -z "$ENV_NAME" ]; then
    echo "[ERROR] No se especificó el nombre del entorno."
    return 1
fi

if [ -z "$BASE_DIR" ]; then
    echo "[ERROR] No se especificó el directorio base."
    return 1
fi

# Verifique si existe el directorio del proyecto
if [ ! -d "$BASE_DIR" ]; then
    echo "[ERROR] Directorio no encontrado: ${BASE_DIR}"
    return 1
fi

# Cambie al directorio del proyecto
cd "$BASE_DIR" || return 1

# Construya la ruta al script de activación
VENV_ROOT="$HOME/penv"
VENV_ACTIVATE="${VENV_ROOT}/${ENV_NAME}/bin/activate"

# Active el entorno
if [ -f "$VENV_ACTIVATE" ]; then
    echo "Activando..."
    # shellcheck source=/dev/null
    source "$VENV_ACTIVATE"
    echo "Activación completa."
else
    echo "[ERROR] No encontrado: ${VENV_ACTIVATE}"
    return 1
fi

# Abra VS Code en el directorio del proyecto (con soporte para WSL)
echo "Iniciando VS Code..."
code .
\end{lstlisting}

Observe varias diferencias respecto de la versión de Windows:
\begin{itemize}
  \item \texttt{-z} comprueba si una cadena está vacía (longitud cero).
  \item \texttt{-d} comprueba si una ruta es un directorio.
  \item \texttt{-f} comprueba si una ruta es un archivo regular.
  \item \texttt{\$HOME} es el equivalente en Linux de
        \texttt{\%USERPROFILE\%}: se expande a su directorio de inicio.
  \item \texttt{return 1} sale del script con un código de error (usamos
        \texttt{return} en lugar de \texttt{exit} por una razón explicada
        en la siguiente sección).
\end{itemize}

\subsubsection{\ProjectName{}.sh}

\begin{lstlisting}[style=linuxStyle, caption={\texttt{\ProjectName.sh}: lanzador del proyecto (Linux/WSL)},escapeinside={(*@}{@*)}]
#!/bin/bash
# (*@\ProjectName@*).sh
# Uso: source ~/(*@\ProjectName@*)/(*@\ProjectName@*).sh
# Inicia el entorno del proyecto (*@\ProjectName@*).

source ~/penv/ActivateEnv.sh (*@\ProjectName@*) "$HOME/(*@\ProjectName@*)"
\end{lstlisting}

\subsection{Hacer ejecutable un script}
\label{sec:chmod}

En Linux, crear un archivo no lo convierte automáticamente en ejecutable.
Debe conceder permiso de ejecución con el comando \texttt{chmod}:

\begin{lstlisting}[style=linuxStyle,escapeinside={(*@}{@*)}]
chmod +x ~/penv/ActivateEnv.sh
chmod +x ~/(*@\ProjectName@*)/(*@\ProjectName@*).sh
\end{lstlisting}

\texttt{chmod} significa \emph{cambiar el modo} (change mode). La bandera
\texttt{+x} añade el permiso de ejecución (\texttt{x}). Los permisos de
Linux se organizan en torno a tres roles (propietario, grupo y otros),
pero para scripts personales, \texttt{chmod +x} es casi siempre todo lo
que necesita.

\subsection{Añadir sus scripts al PATH}

En Windows, añadió \texttt{C:\textbackslash penv} al \texttt{PATH} del
sistema para que \texttt{ActivateEnv.bat} pudiera encontrarse desde
cualquier ubicación. En Linux, hace lo mismo editando el archivo de
configuración de su intérprete de comandos.

\begin{lstlisting}[style=linuxStyle, caption={Añadir \textasciitilde/penv al PATH en \textasciitilde/.bashrc},escapeinside={(*@}{@*)}]
# Abra su archivo bashrc en un editor de texto
nano ~/.bashrc

# Añada esta línea al final:
export PATH="$HOME/penv:$PATH"

# Guarde y recargue la configuración
source ~/.bashrc
\end{lstlisting}

Después de esto, el intérprete de comandos puede encontrar
\texttt{ActivateEnv.sh} desde cualquier directorio.

% ------------------------------------------------------------
\subsection{El enfoque con \texttt{uv}: \texttt{ActivateUVEnv}}
\label{subsec:activateuvenv}
% ------------------------------------------------------------

Los siguientes scripts reflejan a \texttt{ActivateEnv}, pero están dirigidos
a entornos gestionados por \texttt{uv} almacenados en
\texttt{C:\textbackslash uvenv} (Windows) o
\texttt{\textasciitilde/uvenv} (Linux). Almacene todos los entornos de
\texttt{uv} en \texttt{uvenv} en lugar de en \texttt{penv} para mantener
las dos cadenas de herramientas visualmente distintas en el disco.

Dos comportamientos distinguen a \texttt{ActivateUVEnv} de su contraparte
basada en \texttt{pip}:

\begin{enumerate}
  \item \textbf{Creación automática.} Si el entorno con el nombre indicado
        no existe todavía, \texttt{uv venv} lo crea antes de la activación.
        No hay un paso de configuración separado para proyectos nuevos.
  \item \textbf{Sincronización automática.} Si hay un \texttt{pyproject.toml}
        presente en el directorio del proyecto, \texttt{uv sync} se ejecuta
        después de la activación. Las dependencias se instalan por sí
        mismas después de cada \texttt{git pull} que añada nuevos
        requisitos, por lo que no se necesita ningún \texttt{pip install}
        manual.
\end{enumerate}

\subsubsection{ActivateUVEnv.bat (Windows)}

Almacene este archivo en \texttt{C:\textbackslash uvenv} y añada esa
carpeta al \texttt{PATH} del sistema.

\begin{lstlisting}[style=windowsStyle, basicstyle=\footnotesize\ttfamily, caption={\texttt{ActivateUVEnv.bat}: lanzador de entorno uv (Windows)},escapeinside={(*@}{@*)}]
@echo off
setlocal

set "ENV_NAME=%~1"
set "BASE_DIR=%~2"
set "VENV_ROOT=C:\uvenv"
set "VENV_PATH=%VENV_ROOT%\%ENV_NAME%"
set "VENV_ACTIVATE=%VENV_PATH%\Scripts\activate.bat"

:: Valide las entradas
if "%ENV_NAME%"=="" (
    echo [ERROR] Uso: ActivateUVEnv ^<env-name^> ^<project-dir^>
    goto :eof
)
if "%BASE_DIR%"=="" (
    echo [ERROR] Uso: ActivateUVEnv ^<env-name^> ^<project-dir^>
    goto :eof
)
if not exist "%BASE_DIR%" (
    echo [ERROR] Directorio del proyecto no encontrado: %BASE_DIR%
    goto :eof
)

:: Cree el entorno si aún no existe
if not exist "%VENV_ACTIVATE%" (
    echo [INFO] Entorno '%ENV_NAME%' no encontrado. Creándolo con uv...
    uv venv "%VENV_PATH%" --python 3.12
    if errorlevel 1 (
        echo [ERROR] No se pudo crear el entorno.
        goto :eof
    )
)

:: Active
cd /d "%BASE_DIR%"
call "%VENV_ACTIVATE%"

:: Sincronice las dependencias si está presente pyproject.toml
if exist "pyproject.toml" (
    echo [INFO] pyproject.toml encontrado. Sincronizando dependencias...
    uv sync
)

echo.
echo [OK] Entorno '%ENV_NAME%' activo.
echo Iniciando VS Code...
code .

endlocal
\end{lstlisting}

El uso es idéntico en forma al de \texttt{ActivateEnv.bat}:

\begin{lstlisting}[style=windowsStyle,escapeinside={(*@}{@*)}]
ActivateUVEnv lunar "C:\Dropbox\Research\Lunar\python"
ActivateUVEnv forecasting "C:\Projects\forecasting"
\end{lstlisting}

\subsubsection{ActivateUVEnv: función de intérprete de comandos para Linux / WSL}

Añada la siguiente función a \texttt{\textasciitilde/.bashrc}. Al igual que
su contraparte basada en \texttt{pip}, se ejecuta en el intérprete de
comandos actual, de modo que la activación y los cambios de directorio
persisten en su sesión.

\begin{lstlisting}[style=linuxStyle, basicstyle=\footnotesize\ttfamily, caption={Función de intérprete de comandos \texttt{ActivateUVEnv()} para \textasciitilde/.bashrc},escapeinside={(*@}{@*)}]
# -------------------------------------------------------
# ActivateUVEnv: Activa un venv gestionado por uv y abre VS Code
# Uso: ActivateUVEnv <ENV_NAME> <PROJECT_DIR>
# -------------------------------------------------------
ActivateUVEnv() {
    local ENV_NAME="$1"
    local BASE_DIR="$2"
    local VENV_ROOT="$HOME/uvenv"
    local VENV_PATH="${VENV_ROOT}/${ENV_NAME}"
    local VENV_ACTIVATE="${VENV_PATH}/bin/activate"

    if [ -z "$ENV_NAME" ]; then
        echo "[ERROR] Uso: ActivateUVEnv <env-name> <project-dir>"
        return 1
    fi

    if [ -z "$BASE_DIR" ]; then
        echo "[ERROR] Uso: ActivateUVEnv <env-name> <project-dir>"
        return 1
    fi

    if [ ! -d "$BASE_DIR" ]; then
        echo "[ERROR] Directorio del proyecto no encontrado: ${BASE_DIR}"
        return 1
    fi

    # Cree el entorno si aún no existe
    if [ ! -f "$VENV_ACTIVATE" ]; then
        echo "[INFO] Entorno '${ENV_NAME}' no encontrado. Creándolo con uv..."
        uv venv "$VENV_PATH" --python 3.12 || return 1
    fi

    # Active
    cd "$BASE_DIR" || return 1
    # shellcheck source=/dev/null
    source "$VENV_ACTIVATE"

    # Sincronice las dependencias si está presente pyproject.toml
    if [ -f "pyproject.toml" ]; then
        echo "[INFO] pyproject.toml encontrado. Sincronizando dependencias..."
        uv sync
    fi

    echo "[OK] Entorno '${ENV_NAME}' activo."
    echo "Iniciando VS Code..."
    code .
}
\end{lstlisting}

Recargue \texttt{.bashrc} después de guardar:

\begin{lstlisting}[style=linuxStyle,escapeinside={(*@}{@*)}]
source ~/.bashrc
\end{lstlisting}

Luego invóquela desde cualquier directorio:

\begin{lstlisting}[style=linuxStyle,escapeinside={(*@}{@*)}]
ActivateUVEnv lunar "$HOME/Research/Lunar/python"
\end{lstlisting}

% ============================================================
\section{La diferencia crítica: ejecutar frente a cargar con source}
\label{sec:source-vs-execute}
% ============================================================

Esta sección explica uno de los conceptos más importantes, y a la vez más
incomprendidos, de los scripts de intérprete de comandos de Linux,
especialmente para desarrolladores que vienen de Windows.

\subsection{Intérpretes de comandos hijos}

Cada vez que usted ejecuta un script con \texttt{./script.sh} o simplemente
con \texttt{script.sh}, Linux lanza un \textbf{intérprete de comandos hijo}:
un proceso de intérprete de comandos nuevo y separado que hereda una copia
de su entorno, ejecuta el script y después \emph{termina}. Cualquier cambio
que el script realice en el entorno (activar un venv, cambiar de
directorio, establecer una variable) existe únicamente en ese intérprete
hijo. Cuando el hijo termina, esos cambios desaparecen. Su terminal
original no se ve afectado en absoluto.

Por eso ejecutar \texttt{./ActivateEnv.sh} parecería no hacer nada útil:
el venv se activa dentro del hijo, luego el hijo termina, y su terminal
sigue apuntando al Python global.

\subsection{Cargar con source: ejecución en el intérprete de comandos actual}

El comando \texttt{source} (o su abreviatura, un único punto \texttt{.})
ejecuta un script \emph{en el intérprete de comandos actual}: no se crea
ningún hijo. Cada comando se ejecuta como si lo hubiera escrito
directamente en la indicación. Esto significa:

\begin{itemize}
  \item \texttt{source ~/penv/\ProjectName{}/bin/activate} activa el venv
        en \emph{su} terminal.
  \item \texttt{cd} cambia \emph{su} directorio de trabajo.
  \item Las variables establecidas por el script persisten en \emph{su} sesión.
\end{itemize}

\begin{lstlisting}[style=linuxStyle,escapeinside={(*@}{@*)}]
# Estos dos son equivalentes:
source ActivateEnv.sh (*@\ProjectName@*) ~/(*@\ProjectName@*)
. ActivateEnv.sh (*@\ProjectName@*) ~/(*@\ProjectName@*)

# Esto NO funciona para la activación del venv:
./ActivateEnv.sh (*@\ProjectName@*) ~/(*@\ProjectName@*)    # intérprete hijo; los cambios se pierden
\end{lstlisting}

\subsection{El paralelo en Windows}

Windows no tiene exactamente esta distinción de la misma manera. El comando
\texttt{call} en un archivo por lotes ejecuta otro \texttt{.bat} en la
misma sesión de CMD y regresa, razón por la cual la versión de Windows
funciona como se espera. El comportamiento por defecto de Linux de generar
un proceso hijo es en realidad el diseño \emph{más seguro}, ya que evita
que los scripts modifiquen accidentalmente su entorno, pero le exige optar
explícitamente por el comportamiento de sesión compartida mediante
\texttt{source}.

\begin{notebox}
Esta es también la razón por la que \texttt{ActivateEnv.sh} usa
\texttt{return} en lugar de \texttt{exit} para el manejo de errores.
Cuando un script se carga con source, \texttt{exit} cerraría toda su
terminal. \texttt{return} sale únicamente del script, dejando su sesión
intacta.
\end{notebox}

% ============================================================
\section{Nivel 2: funciones de intérprete de comandos, el enfoque elegante de Linux}
\label{sec:level2}
% ============================================================

Cargar scripts con source es poderoso, pero requiere escribir \texttt{source}
o \texttt{.} cada vez. También implica que su script lanzador debe residir
en algún lugar del \texttt{PATH} y que el usuario debe recordar la
invocación exacta. Linux ofrece una solución más elegante: las
\textbf{funciones de intérprete de comandos}.

\subsection{¿Qué es una función de intérprete de comandos?}

Una función de intérprete de comandos es como un script, pero se define
directamente dentro de su archivo de configuración del intérprete de
comandos (\texttt{\textasciitilde/.bashrc}). Debido a que forma parte del
propio intérprete (y no es un archivo separado), se ejecuta automáticamente
en la sesión actual sin necesidad de \texttt{source}. Usted la invoca
simplemente escribiendo su nombre.

\subsection{Definición de la función ActivateEnv}

Añada lo siguiente al final de \texttt{\textasciitilde/.bashrc}:

\begin{lstlisting}[style=linuxStyle, basicstyle=\footnotesize\ttfamily, caption={Versión como función de intérprete de comandos de \texttt{ActivateEnv} en \textasciitilde/.bashrc},escapeinside={(*@}{@*)}]
# -------------------------------------------------------
# ActivateEnv: Activa un venv de Python y abre VS Code
# Uso: ActivateEnv <ENV_NAME> <PROJECT_DIR>
# -------------------------------------------------------
ActivateEnv() {
    local ENV_NAME="$1"
    local BASE_DIR="$2"

    echo "Iniciando la activación del entorno '${ENV_NAME}'..."

    # Valide la entrada
    if [ -z "$ENV_NAME" ]; then
        echo "[ERROR] No se especificó el nombre del entorno."
        return 1
    fi

    if [ -z "$BASE_DIR" ]; then
        echo "[ERROR] No se especificó el directorio base."
        return 1
    fi

    if [ ! -d "$BASE_DIR" ]; then
        echo "[ERROR] Directorio no encontrado: ${BASE_DIR}"
        return 1
    fi

    local VENV_ACTIVATE="$HOME/penv/${ENV_NAME}/bin/activate"

    if [ ! -f "$VENV_ACTIVATE" ]; then
        echo "[ERROR] Entorno virtual no encontrado: ${VENV_ACTIVATE}"
        return 1
    fi

    # Cambie de directorio, active y abra VS Code
    cd "$BASE_DIR" || return 1
    source "$VENV_ACTIVATE"
    echo "Entorno '${ENV_NAME}' activado."
    code .
}
\end{lstlisting}

La palabra clave \texttt{local} restringe las variables al ámbito de la
función, de modo que no contaminan la sesión del intérprete de comandos.

\subsection{Definición de la función \ProjectName{}}

Ahora añada una función de una línea que invoque a \texttt{ActivateEnv}
para el proyecto \ProjectName{}:

\begin{lstlisting}[style=linuxStyle, caption={Función específica del proyecto para \ProjectName{} en \textasciitilde/.bashrc},escapeinside={(*@}{@*)}]
# -------------------------------------------------------
# (*@\ProjectName@*): Inicia el entorno del proyecto (*@\ProjectName@*)
# Uso: (*@\ProjectName@*)
# -------------------------------------------------------
(*@\ProjectName@*)() {
    ActivateEnv "(*@\ProjectName@*)" "$HOME/(*@\ProjectName@*)"
}
\end{lstlisting}

\subsection{Activación de las nuevas funciones}

Después de editar \texttt{\textasciitilde/.bashrc}, recárguelo:

\begin{lstlisting}[style=linuxStyle,escapeinside={(*@}{@*)}]
source ~/.bashrc
\end{lstlisting}

Ahora, desde \emph{cualquier directorio}, simplemente escriba:

\begin{lstlisting}[style=linuxStyle,escapeinside={(*@}{@*)}]
(*@\ProjectName@*)
\end{lstlisting}

El intérprete de comandos cambiará al directorio de su proyecto \ProjectName{},
activará el entorno virtual y abrirá VS Code, todo en una sola palabra y
sin que se requiera el prefijo \texttt{source}.

\subsection{Por qué Windows no tiene un equivalente limpio}

Los archivos por lotes de Windows siempre se ejecutan como procesos hijos;
no existe el comando \texttt{source} en CMD. PowerShell es lo más cercano:
puede definir funciones en su perfil de PowerShell (\texttt{\$PROFILE}) y
se comportan de manera similar a las funciones de Bash. Sin embargo, el
perfil de PowerShell está deshabilitado por defecto por razones de
seguridad y requiere cambiar la política de ejecución. Para los usuarios
de Windows que desean esta experiencia de un solo comando de forma nativa,
configurar los perfiles de PowerShell es el camino a seguir, pero implica
más sobrecarga administrativa que editar
\texttt{\textasciitilde/.bashrc}. Esta es una de las razones prácticas
por las que muchos desarrolladores en Windows adoptan WSL.

% ============================================================
\section{Juntando todas las piezas}
\label{sec:together}
% ============================================================

La Figura~\ref{fig:workflow} resume el flujo de trabajo completo en cada
plataforma.

\begin{table}[h]
\centering
\caption{Comparación completa del flujo de trabajo: Windows frente a Linux/WSL}
\label{fig:workflow}
\begin{tabular}{@{}p{0.15\textwidth}p{0.38\textwidth}p{0.38\textwidth}@{}}
\toprule
\textbf{Paso} & \textbf{Windows} & \textbf{Linux / WSL} \\
\midrule
Crear venv
  & \texttt{python -m venv C:\textbackslash penv\textbackslash \ProjectName}
  & \texttt{python3 -m venv \textasciitilde/penv/\ProjectName} \\[4pt]
Lanzador reutilizable
  & \texttt{ActivateEnv.bat} en \texttt{C:\textbackslash penv}
  & Función \texttt{ActivateEnv()} en \texttt{\textasciitilde/.bashrc} \\[4pt]
Lanzador de proyecto
  & \texttt{\ProjectName{}.bat} en la carpeta del proyecto
  & Función \texttt{\ProjectName{}()} en \texttt{\textasciitilde/.bashrc} \\[4pt]
Hacerlo disponible globalmente
  & Añadir \texttt{C:\textbackslash penv} al PATH del sistema
  & Las funciones en \texttt{.bashrc} siempre están disponibles \\[4pt]
Iniciar el proyecto
  & Escribir \texttt{\ProjectName{}} en cualquier ventana de CMD
  & Escribir \texttt{\ProjectName{}} en cualquier terminal \\[4pt]
Resultado
  & venv activo, VS Code abre en la carpeta del proyecto
  & venv activo, VS Code abre en la carpeta del proyecto (modo WSL) \\
\bottomrule
\end{tabular}
\end{table}

\subsection{La experiencia completa en WSL}

Cuando usted escribe \texttt{\ProjectName{}} en su terminal de WSL:

\begin{enumerate}
  \item Bash busca \texttt{\ProjectName{}} en su lista de funciones conocidas.
  \item \texttt{\ProjectName{}()} invoca a \texttt{ActivateEnv "\ProjectName{}" "\$HOME/\ProjectName{}"}.
  \item \texttt{ActivateEnv()} cambia al directorio del proyecto.
  \item El entorno virtual se activa en su intérprete de comandos actual.
  \item \texttt{code .} inicia VS Code con la extensión de WSL, conectando
        su interfaz de VS Code en Windows con su backend de Linux.
  \item Ahora está trabajando en un entorno de desarrollo completamente
        configurado con una sola palabra.
\end{enumerate}

% ============================================================
\section{Hoja de referencia rápida}
\label{sec:cheatsheet}
% ============================================================

\begin{table}[h]
\centering
\caption{Hoja de referencia de scripts de intérprete de comandos y entornos virtuales}
\label{tab:cheatsheet}
\begin{tabular}{@{}p{0.35\textwidth}p{0.28\textwidth}p{0.28\textwidth}@{}}
\toprule
\textbf{Tarea} & \textbf{Windows (CMD)} & \textbf{Linux / WSL (Bash)} \\
\midrule
Crear entorno virtual
  & \texttt{python -m venv C:\textbackslash penv\textbackslash NAME}
  & \texttt{python3 -m venv \textasciitilde/penv/NAME} \\[4pt]
Activar entorno
  & \texttt{C:\textbackslash penv\textbackslash NAME\textbackslash Scripts\textbackslash activate}
  & \texttt{source \textasciitilde/penv/NAME/bin/activate} \\[4pt]
Desactivar entorno
  & \texttt{deactivate}
  & \texttt{deactivate} \\[4pt]
Instalar un paquete
  & \texttt{pip install pkg}
  & \texttt{pip install pkg} \\[4pt]
Guardar lista de paquetes
  & \texttt{pip freeze > req.txt}
  & \texttt{pip freeze > req.txt} \\[4pt]
Restaurar desde la lista
  & \texttt{pip install -r req.txt}
  & \texttt{pip install -r req.txt} \\[4pt]
Hacer ejecutable un script
  & (no es necesario; la extensión lo gestiona)
  & \texttt{chmod +x script.sh} \\[4pt]
Ejecutar un script
  & \texttt{script.bat}
  & \texttt{./script.sh} \\[4pt]
Cargar con source un script
  & \texttt{call script.bat}
  & \texttt{source script.sh} \\[4pt]
Editar configuración del intérprete
  & (sin equivalente en CMD)
  & \texttt{nano \textasciitilde/.bashrc} \\[4pt]
Recargar configuración del intérprete
  & (reiniciar CMD)
  & \texttt{source \textasciitilde/.bashrc} \\[4pt]
Abrir VS Code (WSL)
  & \texttt{code .}
  & \texttt{code .} \\[4pt]
Parámetro 1 del script
  & \texttt{\%\textasciitilde1}
  & \texttt{\$1} \\[4pt]
Comprobar si una variable está vacía
  & \texttt{if "\%VAR\%"==""}
  & \texttt{if [ -z "\$VAR" ]} \\[4pt]
Comprobar si existe un directorio
  & \texttt{if exist "DIR\textbackslash"}
  & \texttt{if [ -d "DIR" ]} \\[4pt]
Salir de un script o función
  & \texttt{goto :eof}
  & \texttt{return} (en función o script cargado con source) \\
\bottomrule
\end{tabular}
\end{table}

\begin{table}[h]
\centering
\caption{Referencia rápida de \texttt{uv}: equivalencias con \texttt{pip} y \texttt{venv}}
\label{tab:uv-cheatsheet}
\begin{tabular}{@{}p{0.32\textwidth}p{0.28\textwidth}p{0.28\textwidth}@{}}
\toprule
\textbf{Tarea} & \textbf{pip / venv} & \textbf{uv} \\
\midrule
Instalar \texttt{uv}
  & (no aplicable)
  & Véase la Sección~\ref{sec:uv} \\[4pt]
Crear entorno
  & \texttt{python -m venv PATH}
  & \texttt{uv venv PATH -{}-python 3.12} \\[4pt]
Activar entorno
  & \texttt{source PATH/bin/activate}
  & \texttt{source PATH/bin/activate} \\[4pt]
Instalar un paquete
  & \texttt{pip install pkg}
  & \texttt{uv pip install pkg} \\[4pt]
Instalar desde \texttt{pyproject.toml}
  & \texttt{pip install -e .}
  & \texttt{uv sync} \\[4pt]
Añadir una dependencia
  & Editar \texttt{requirements.txt} manualmente
  & \texttt{uv add pkg} \\[4pt]
Eliminar una dependencia
  & Editar \texttt{requirements.txt} manualmente
  & \texttt{uv remove pkg} \\[4pt]
Generar archivo de bloqueo
  & \texttt{pip freeze > req.txt}
  & \texttt{uv lock} (o implícito en \texttt{uv sync}) \\[4pt]
Instalar desde archivo de bloqueo
  & \texttt{pip install -r req.txt}
  & \texttt{uv sync} \\[4pt]
Script lanzador (Windows)
  & \texttt{ActivateEnv.bat}
  & \texttt{ActivateUVEnv.bat} \\[4pt]
Función lanzadora (Linux)
  & \texttt{ActivateEnv()} en \texttt{.bashrc}
  & \texttt{ActivateUVEnv()} en \texttt{.bashrc} \\[4pt]
Raíz del entorno
  & \texttt{C:\textbackslash penv} / \texttt{\textasciitilde/penv}
  & \texttt{C:\textbackslash uvenv} / \texttt{\textasciitilde/uvenv} \\
\bottomrule
\end{tabular}
\end{table}

% ============================================================
\section{Resumen}
\label{sec:summary}
% ============================================================

Este capítulo cubrió las herramientas y técnicas para gestionar entornos de
Python y automatizar la configuración de proyectos en Windows y Linux:

\begin{itemize}
  \item \textbf{Entornos virtuales de Python}: copias aisladas de Python,
        una por proyecto, que previenen conflictos de dependencias entre
        proyectos. Cada entorno tiene su propio conjunto de paquetes
        instalados, y eliminar uno no tiene efecto sobre ningún otro.
  \item \textbf{La cadena estándar de herramientas \texttt{pip} y
        \texttt{venv}}: las herramientas respaldadas por la fundación e
        incluidas con cada instalación de Python para crear entornos,
        instalar paquetes y registrar dependencias en
        \texttt{requirements.txt}. Conocer esta cadena de herramientas
        sigue siendo una habilidad esencial independientemente de qué
        otras herramientas adopte.
  \item \textbf{\texttt{uv}}: una alternativa moderna y de alto rendimiento
        desarrollada por Astral que ofrece instalaciones drásticamente más
        rápidas, una caché global de paquetes, archivos de bloqueo
        reproducibles con hashes criptográficos y soporte nativo para
        \texttt{pyproject.toml}. Debido a que \texttt{uv} sincroniza
        automáticamente las dependencias desde un archivo de bloqueo
        mediante \texttt{uv sync}, la reproducción del entorno se vuelve
        un único comando en lugar de un proceso manual.
  \item \textbf{Scripting de Nivel~1}: escribir scripts lanzadores
        reutilizables (\texttt{ActivateEnv.bat} y \texttt{ActivateEnv.sh}
        para entornos gestionados por \texttt{pip};
        \texttt{ActivateUVEnv.bat} y la función de intérprete de comandos
        \texttt{ActivateUVEnv()} para entornos gestionados por
        \texttt{uv}) y scripts de una línea específicos del proyecto
        (\texttt{\ProjectName{}.bat}, \texttt{\ProjectName{}.sh}) que
        activan el entorno correcto, cambian al directorio del proyecto
        e inician VS Code en una sola invocación.
  \item \textbf{Scripting de Nivel~2}: funciones de intérprete de comandos
        definidas en \texttt{\textasciitilde/.bashrc} que proporcionan un
        comando de inicio de una sola palabra desde cualquier directorio,
        sin necesidad del prefijo \texttt{source}.
  \item \textbf{La distinción entre \texttt{source} y ejecutar}: en Linux,
        ejecutar un script normalmente genera un intérprete de comandos
        hijo cuyos cambios de entorno se desvanecen cuando termina. El
        comando \texttt{source} (o su abreviatura \texttt{.}) ejecuta el
        script en el intérprete de comandos actual, razón por la cual la
        activación del entorno virtual requiere cargarlo con source. Esta
        misma distinción explica por qué \texttt{ActivateEnv.sh} usa
        \texttt{return} en lugar de \texttt{exit} para el manejo de
        errores.
\end{itemize}

El comando \texttt{\ProjectName{}} que construyó en este capítulo es un
ejemplo pequeño pero significativo de un principio más amplio: la línea
de comandos le permite codificar su flujo de trabajo en scripts
repetibles, compartibles y mejorables. Cada paso tedioso de configuración
que elimine hoy es carga cognitiva liberada para el trabajo que realmente
importa.

